% Unidad R571: EGA I, capítulo 0, §1.0, (1.0.1)--(1.0.3).
% Traducción: OpenAI Codex, GPT-6.1 Sol, esfuerzo Ultra.
\clearpage
\oldpage[0\textsubscript{I}]{11}
\begin{center}
{\large CAPÍTULO 0\par}
\vspace{1.5em}
{\Large\bfseries PRELIMINARES\par}
\end{center}

\section{Anillos de fracciones}
\label{section:0.1-fr}

\setcounter{subsection}{-1}
\subsection{Anillos y álgebras}
\label{subsection:0.1.0-fr}

\begin{env}[1.0.1]
\label{0.1.0.1-fr}
Todos los anillos considerados en este tratado tendrán un \emph{elemento unidad};
todos los módulos sobre un anillo de este tipo se supondrán \emph{unitarios}; los homomorfismos de anillos siempre se supondrán \emph{que llevan el elemento unidad al elemento unidad}; salvo indicación expresa de lo contrario, se supondrá que un subanillo de un anillo $A$ \emph{contiene el elemento unidad de $A$}.
Consideraremos principalmente anillos \emph{conmutativos}, y cuando hablemos de un anillo sin más precisiones, se entenderá que se trata de un anillo conmutativo.
Si $A$ es un anillo no necesariamente conmutativo, por $A$-módulo entenderemos siempre un módulo \emph{a izquierda}, salvo indicación expresa de lo contrario.
\end{env}

\begin{env}[1.0.2]
\label{0.1.0.2-fr}
Sean $A$, $B$ dos anillos no necesariamente conmutativos, y $\varphi:A\to B$ un homomorfismo.
Todo $B$-módulo a izquierda (resp. a derecha) $M$ puede dotarse de una estructura de $A$-módulo a izquierda (resp. a derecha) poniendo $a.m=\varphi(a).m$ (resp. $m.a=m.\varphi(a)$); cuando sea necesario distinguir en $M$ las estructuras de $A$-módulo y de $B$-módulo, denotaremos por $M_{[\varphi]}$ el $A$-módulo a izquierda (resp. a derecha) así definido.
Si $L$ es un $A$-módulo, un homomorfismo $u:L\to M_{[\varphi]}$ es, por tanto, un homomorfismo de grupos conmutativos tal que $u(a.x)=\varphi(a).u(x)$ para $a\in A$, $x\in L$; también se dirá que es un \emph{$\varphi$-homomorfismo} $L\to M$, y que el par $(\varphi,u)$ (o, por abuso de lenguaje, $u$) es un \emph{dihomomorfismo} de $(A,L)$ a $(B,M)$.
Los pares $(A,L)$ formados por un anillo $A$ y un $A$-módulo $L$ forman, por tanto, una \emph{categoría} cuyos morfismos son los dihomomorfismos.
\end{env}

\begin{env}[1.0.3]
\label{0.1.0.3-fr}
Bajo las hipótesis de (1.0.2), si $\mathfrak{J}$ es un ideal a izquierda (resp. a derecha) de $A$, denotaremos por $B\mathfrak{J}$ (resp. $\mathfrak{J}B$) el ideal a izquierda (resp. a derecha) $B\varphi(\mathfrak{J})$ (resp. $\varphi(\mathfrak{J})B$) de $B$ generado por $\varphi(\mathfrak{J})$; es también la imagen del homomorfismo canónico $B\otimes_A\mathfrak{J}\to B$ (resp. $\mathfrak{J}\otimes_A B\to B$) de $B$-módulos a izquierda (resp. a derecha).
\end{env}
